\documentclass[12pt,a4paper]{{article}}

% --------------------------------------------------
% IDIOMA Y TIPOGRAFÍA
% --------------------------------------------------

\usepackage{fontspec}
\usepackage{polyglossia}

\setdefaultlanguage{spanish}

% Fuente Sans Serif
\setsansfont{Helvetica}
\setmainfont{Helvetica}
\setmonofont{CMU Typewriter Text}

% --------------------------------------------------
% FORMATO GENERAL
% --------------------------------------------------

% Márgenes de 2 cm
\usepackage[
    a4paper,
    top=2cm,
    bottom=2cm,
    left=2cm,
    right=2cm
]{geometry}

% Interlineado 1,5
\usepackage{setspace}
\onehalfspacing

% Incluye imágenes
\usepackage{graphicx}
\usepackage{tcolorbox} % Es para cuadros de colores
\tcbuselibrary{breakable, skins} % Para poder partir un cuadro
\usepackage{subcaption}

% Math package
\usepackage{amsmath}
\usepackage{amssymb}
\usepackage{amsthm}
% Math graphic
\usepackage{pgfplots}
\usepackage{tikz}
\usetikzlibrary{matrix, positioning, calc}
\usepgfplotslibrary{fillbetween}
\usetikzlibrary{positioning, arrows.meta}
\pgfplotsset{compat=1.18}
\usepackage{morefloats}
\usepackage{booktabs}

% Codigo
\usepackage{listings}
\tcbuselibrary{listings} % Es lo que permite usar t colorbox dentro de listings

% Graficos de Matplot
\newcommand{\mathdefault}[1][]{#1}

% Referenciar indice
\usepackage{xcolor}
\usepackage{sectsty}

\usepackage[colorlinks=true,
    linkcolor=black,
    urlcolor=black,
    citecolor=black,
    hypertexnames=false]{hyperref}

% Citas a bibliografia o fuentes
\usepackage[
    backend=biber,  % Motor moderno para procesar las citas
    language=spanish % Para que diga "y" en vez de "and" entre autores
]{biblatex}
\addbibresource{bibliografia.bib}
\usepackage{url}
\usepackage{hyperref}
\hypersetup{
    colorlinks=true,
    linkcolor=black,   % Cambia a black para que el índice no sea azul
    urlcolor=blue,     
    citecolor=blue,
    allbordercolors={white} % Quita recuadros si existieran
}
\setcounter{biburllcpenalty}{7000}
\setcounter{biburlucpenalty}{8000}

\usepackage{titlesec}

\titleformat{\section}
{\color{red}\normalfont\Large\bfseries}
{\thesection}{1em}{}

\titleformat{\subsection}
{\color{red}\normalfont\large\bfseries}
{\thesubsection}{1em}{}

\titleformat{\subsubsection}
{\color{red}\normalfont\normalsize\bfseries}
{\thesubsubsection}{1em}{}

% Entornos
%=============================
% ENTORNOS PERSONALIZADOS

% Cuadro ejemplo
\newcounter{example}[section]
\renewcommand{\theexample}{\thesection.\arabic{example}}

\newtcolorbox{ejemplo}[1][]{
    enhanced,
    breakable,
    colback=gray!12,
    boxrule=0pt,
    frame hidden,
    borderline west={2pt}{0pt}{black},
    sharp corners,
    left=8pt,
    right=8pt,
    top=6pt,
    bottom=6pt,
    fonttitle=\bfseries,
    title={Ejemplo~\theexample},
    coltitle=black,
    code={\refstepcounter{example}},
    #1
}

% Notas
\newtcolorbox{nota}{
  enhanced,
  breakable,
  colback=gray!6,
  boxrule=0pt,
  frame hidden,
  borderline west={3pt}{0pt}{gray!60},
  sharp corners,
  left=8pt,
  right=8pt,
  top=6pt,
  bottom=6pt,
  fontupper=\small
}

% Ejercicios 
\newtcolorbox{ejercicio}[1][]{
  enhanced,
  breakable,
  colback=white,
  colframe=black!70,      % Color del marco y del título
  colbacktitle=black!70, % Fondo del título
  coltitle=white,        % Texto del título en blanco
  fonttitle=\bfseries\small,
  title={EJERCICIO #1},   % El título es "EJERCICIO" + lo que tú escribas
  detach title,
  % Diseño del encabezado
  before upper={\textsf{\textbf{\textcolor{black!70}{EJERCICIO #1}}}\smallskip\hrule\medskip\par},
  boxrule=0pt,
  frame hidden,
  borderline north={1.5pt}{0pt}{black!70}, % Línea superior gruesa
  borderline south={0.5pt}{0pt}{black!20}, % Línea inferior sutil
  sharp corners,
  left=0pt,   % Alineado a la izquierda sin sangría
  right=0pt,
  top=5pt,
  bottom=5pt,
  fontupper=\small
}


% Teorema
\theoremstyle{definition}
\newtheorem{theorem}{Teorema}%[section]


% ----- CODIGO ----- %

% 1. Definimos los colores para Python
\definecolor{codegreen}{rgb}{0,0.6,0}
\definecolor{codepurple}{rgb}{0.58,0,0.82}
\definecolor{backgray}{rgb}{0.95,0.95,0.96}
\definecolor{pyblue}{rgb}{0.1, 0.4, 0.6}

% 2. Creamos el entorno personalizado llamado "pycode"
\newtcblisting{pycode}[1]{%
    breakable,
    enhanced,
    arc=2pt,
    outer arc=2pt,
    boxrule=0.5pt,
    colback=backgray,
    colframe=pyblue!70!black,
    left=5mm,
    overlay={
        \begin{tcbclipinterior}
            \fill[gray!20] (frame.south west) rectangle ([xshift=5mm]frame.north west);
        \end{tcbclipinterior}
    },
    title=#1, % El título que le pases al entorno
    fonttitle=\bfseries\small\ttfamily,
    listing only,
    listing options={
        language=Python,
        basicstyle=\small\ttfamily,
        keywordstyle=\color{blue}\bfseries,
        stringstyle=\color{codepurple},
        commentstyle=\color{codegreen}\itshape,
        numbers=left,
        numberstyle=\tiny\color{gray},
        stepnumber=1,
        numbersep=8pt,
        showstringspaces=false,
        breaklines=true,
        tabsize=4
    }
}

%=============================
\title{\color{red}Unidad 10 - Grupo 17; Detección de Fraude}
\author{\itshape Gino Anci, Franco Ricci, Sánchez Facundo, Cossio Pablo}
\date{\small\today}

\begin{document}
    
\begin{titlepage}
    \newgeometry{a4paper, margin=2.5cm}
    \thispagestyle{empty}
    \sffamily
    \centering

    \vspace*{1.2cm}

    \begin{center}
    \begin{tikzpicture}
        \node[red, font=\Large\sffamily,scale=2] at(6,0) {\textbf{UTN}};
        \draw[red, very thick] (7.6,0.6) -- (7.6,-0.6);
        \node[right=7pt, font=\large\sffamily, text width=6cm] at (7.6,0) {Universidad\\Tecnológica Nacional};
    \end{tikzpicture}
    \end{center}

    \vfill

    {\large\bfseries Diplomatura en Ciencia de Datos y Análisis Avanzado\par}

    \vspace{0.35cm}

    {\normalsize Unidad 10 --- Proyecto Final Integrador\par}

    \vspace{0.9cm}

    {\color{red}\rule{1.5cm}{1.5pt}\par}

    \vspace{0.9cm}

    % Título principal
    {\bfseries\fontsize{21}{28}\selectfont
        Detección de fraude en transacciones bancarias mediante aprendizaje automático\par}

    \vspace{1.8cm}

    {\color{black}\large GRUPO 17\par}

    \vspace*{0.6cm}

    % Integrantes
    {\small INTEGRANTES\par}

    \vspace{0.35cm}

    {\normalsize
        Gino Anci - Franco Ricci - Facundo Sánchez - Pablo Cossio\par}
    \vfill

    % Año
    {\small2026\par}

    \vspace*{0.3cm}

    \restoregeometry
\end{titlepage}


    \tableofcontents    

\section{Resumen ejecutivo}

El presente proyecto aborda la detección de fraude en transacciones financieras mediante técnicas de aprendizaje automático supervisado. A partir del conjunto de datos sintético \textit{BankSim}, se desarrollaron y compararon modelos capaces de asignar un nivel de riesgo a las operaciones, con el objetivo de contribuir a la identificación de clientes potencialmente comprometidos.

El trabajo se orienta a un escenario de prevención de fraude en una entidad financiera, donde la capacidad de revisar alertas es limitada. Se utilizó como indicador operativo la métrica \textit{Card Precision at k} (CP@$k$), que expresa qué proporción de los $k$ clientes seleccionados diariamente presenta fraude. Se estudiaron capacidades de 100 y 20 revisiones por día, complementando el análisis mediante \textit{Average Precision} (AP) y ROC-AUC, de acuerdo con la metodología propuesta en \cite{leborgne2022fraud}.

Se compararon una regresión logística como modelo de referencia, \textit{Random Forest}, XGBoost y una red neuronal \textit{Feed-Forward}. El mejor resultado de validación para CP@20 correspondió a la red neuronal con feature engineering (0,3000). Este modelo se evaluó finalmente en un conjunto temporal de prueba independiente y obtuvo CP@20 de 0,2217, AP de 0,8902 y ROC-AUC de 0,9993.

En la prueba, el procedimiento permitió identificar 102 de los 104 clientes con fraude presentes en el conjunto evaluable, considerando un máximo de 20 revisiones diarias durante 23 días. Este resultado describe el desempeño sobre datos simulados y no debe interpretarse como una estimación directa del ahorro económico ni como garantía de resultados equivalentes en una entidad real.
\section{Definición del problema y relevancia}

El fraude en los medios de pago representa un problema para las entidades financieras debido a las operaciones no autorizadas, las posibles pérdidas económicas y la necesidad de destinar recursos a su investigación. Su detección presenta dificultades asociadas al desbalance entre transacciones legítimas y fraudulentas, a la evolución de los patrones de fraude y a las restricciones operativas de los equipos de prevención \cite{leborgne2022fraud}.

El proyecto plantea un caso de negocio hipotético: una entidad financiera dispone de un conjunto de transacciones y debe decidir cuáles revisar primero. No resulta suficiente clasificar correctamente un porcentaje elevado de operaciones, dado que la mayoría es legítima. El problema consiste en ordenar las alertas de manera que las revisiones disponibles se concentren sobre clientes realmente comprometidos.

\subsection{Objetivo e hipótesis}
El objetivo es desarrollar y evaluar modelos de clasificación que permitan detectar operaciones fraudulentas, buscando maximizar la cantidad de clientes comprometidos identificados y limitar las revisiones de operaciones legítimas. De manera medible, se propone superar el desempeño de una regresión logística de referencia, dando prioridad a CP@20 ante un escenario de capacidad operativa restringida.

Se plantea como hipótesis que las características de las transacciones y los antecedentes recientes de clientes, comercios y categorías permiten reconocer patrones asociados al fraude. Para contrastarla, se compararon modelos entrenados con las variables originales disponibles y modelos que incorporan variables históricas realizadas con feature engineering.

\subsection{Alcance y restricciones}
El proyecto no contempla el despliegue de un sistema productivo. Se trabaja con información sintética, por lo que los resultados sirven para estudiar y comparar alternativas metodológicas, pero no permiten validar su rendimiento sobre transacciones bancarias reales. Además, el dataset identifica a cada \texttt{customer} sin proporcionar un identificador independiente de tarjeta. Por este motivo, se adopta el cliente como unidad equivalente a la tarjeta o cuenta comprometida al calcular CP@$k$.

Se considera un retraso de siete días para disponer de las etiquetas de fraude. Esta restricción se incorpora tanto a la generación de antecedentes de riesgo como a la simulación de clientes previamente identificados.
\section{Datos y metodología}

El desarrollo se organizó comprendiendo el problema de negocio, comprendiendo y preparando los datos, modelando y evaluando. La fase de despliegue se dejó fuera del alcance, dado que el objetivo es comparar alternativas sobre una base reproducible.

\subsection{Fuente y relevamiento de datos}
El conjunto de datos utilizado es \textit{BankSim}, disponible públicamente en Kaggle \cite{dataset}. Fue desarrollado por López-Rojas y Axelsson \cite{inproceedings} como un simulador de pagos bancarios que permite generar transacciones sintéticas para investigaciones sobre detección de fraude.

El archivo analizado contiene 594.643 transacciones y diez columnas, con operaciones distribuidas entre los días 0 y 179 de la simulación. La variable objetivo es \texttt{fraud}, que indica si una transacción fue etiquetada como fraudulenta.

\begin{table}[!htbp]
\centering
\small
\caption{Variables originales del conjunto BankSim}
\label{tab:variables}
\begin{tabular}{@{}p{2.65cm}p{6.0cm}p{3.3cm}@{}}
\toprule
\textbf{Variable} & \textbf{Descripción} & \textbf{Tipo} \\
\midrule
\texttt{step} & Día de la simulación & Temporal \\
\texttt{customer} & Identificador del cliente & Identificador \\
\texttt{age} & Rango de edad del cliente & Categórica \\
\texttt{gender} & Género o tipo de cliente & Categórica \\
\texttt{zipcodeOri} & Código postal del cliente & Categórica \\
\texttt{merchant} & Identificador del comercio & Identificador \\
\texttt{zipMerchant} & Código postal del comercio & Categórica \\
\texttt{category} & Categoría de compra o comercio & Categórica \\
\texttt{amount} & Importe de la transacción & Numérica \\
\texttt{fraud} & Etiqueta de fraude (0 o 1) & Objetivo binario \\
\bottomrule
\end{tabular}
\par\smallskip
\end{table}

\subsection{Exploración inicial y calidad}
El análisis de datos exploratorio (EDA) comenzó con la verificación de los tipos de variables, las dimensiones del conjunto y la presencia de valores faltantes. No se observaron valores nulos en las columnas conservadas para el modelado. Los textos se normalizaron eliminando las comillas simples que delimitaban sus valores.

Se eliminaron \texttt{age} y \texttt{gender}, dado que los autores de BankSim aclaran que el comportamiento de consumo de los agentes no se programó en función de estas características \cite{inproceedings}. También se descartaron \texttt{zipcodeOri} y \texttt{zipMerchant}, que en el anteproyecto se identificaron como variables sin variación. Tras esta etapa quedaron seis columnas originales, incluida la variable objetivo.

Como aproximación inicial a la relación con el fraude, se obtuvo una correlación de Pearson de 0,490 entre \texttt{amount} y \texttt{fraud}. Para las categorías, analizadas como indicadores binarios, se observaron asociaciones positivas en \texttt{es\_sportsandtoys} (0,364), \texttt{es\_travel} (0,250) y \texttt{es\_leisure} (0,248), y negativa en \texttt{es\_transportation} ($-0,263$). Estas relaciones son descriptivas del dataset y no implican causalidad. Asimismo, no se encontraron valores atípicos que pertenezcan a errores, por lo que se decidio conservar la mayor cantidad de información.

\subsection{Ingeniería de variables}
Siguiendo los lineamientos del manual de detección de fraude \cite{leborgne2022fraud}, se construyeron dieciocho variables adicionales para incorporar información histórica en ventanas de 1, 7 y 30 días. El procedimiento consideró tres niveles de agregación:

\begin{table}[!htbp]
\centering
\small
\caption{Variables generadas mediante \textit{feature engineering}}
\label{tab:features}
\begin{tabular}{@{}p{2.65cm}p{6.55cm}p{6cm}@{}}
\toprule
\textbf{Entidad} & \textbf{Variables} & \textbf{Información representada} \\
\midrule
Cliente & \texttt{tx\_count\_last\_nd}; \texttt{avg\_amount\_last\_nd} & Cantidad y monto promedio de transacciones recientes. \\
Comercio & \texttt{merchant\_tx\_delayed\_nd}; \texttt{merchant\_risk\_nd} & Actividad y proporción de fraudes en antecedentes disponibles. \\
Categoría & \texttt{category\_tx\_delayed\_nd}; \texttt{category\_risk\_nd} & Actividad y proporción de fraude de la categoría. \\
\bottomrule
\end{tabular}
\par\smallskip
{\footnotesize En los nombres, \texttt{n} toma los valores 1, 7 y 30.}
\end{table}

Para calcular el riesgo histórico de un comercio o categoría se utilizó la relación entre fraudes observados y transacciones registradas en una ventana anterior:
\begin{equation}
\label{eq:riesgo}
\operatorname{Riesgo}_{n} = \frac{\text{Fraudes históricos en }n\text{ días}}{\text{Transacciones históricas en }n\text{ días}}.
\end{equation}

Las ventanas de riesgo excluyen el día que se está evaluando y los siete días previos: para una transacción del día $t$, los antecedentes utilizados llegan hasta $t-8$. En ausencia de transacciones en la ventana, el indicador de riesgo se definió como cero. Las ventanas de cliente, en cambio, incluyen la operación actual; esta particularidad debe tenerse en cuenta si el procedimiento se adapta a un entorno de detección en tiempo real con marcas temporales más precisas.

La idea de construir indicadores de importe anómalo respecto del comportamiento individual se consideró, pero no se incorporó a los experimentos finales. Después de la ingeniería de variables, el conjunto contó con 24 columnas antes de aplicar codificación categórica.

\subsection{División temporal y filtrado de clientes}
La separación se realizó por días, evitando mezclar observaciones anteriores y posteriores entre entrenamiento y evaluación. Se dejaron dos intervalos de siete días que representan el retraso de disponibilidad de las etiquetas.

\begin{table}[!htbp]
\centering
\small
\caption{Partición temporal y efecto del filtrado}
\label{tab:splits}
\begin{tabular}{@{}lrrrr@{}}
\toprule
\textbf{Período} & \textbf{Días} & \textbf{Transacciones} & \textbf{Fraudes} & \textbf{Tras filtrado} \\
\midrule
Entrenamiento & 0--119 & 374.914 & 4.800 & 374.914 \\
Espera & 120--126 & 24.986 & 280 & No evaluado \\
Validación & 127--149 & 83.437 & 920 & 60.655 \\
Espera & 150--156 & --- & --- & No evaluado \\
Prueba & 157--179 & 85.380 & 920 & 58.305 \\
\bottomrule
\end{tabular}
\par\smallskip
{\footnotesize Las columnas de transacciones y fraudes corresponden a los periodos sin filtrar.}
\end{table}

Además de respetar el orden temporal, se eliminaron de validación y prueba las operaciones de clientes cuyo fraude ya debía ser conocido al momento de emitir la alerta. El criterio supone que \textit{una entidad financiera dejaría de investigar una cuenta que ya fue identificada y bloqueada}. Se incorporaron progresivamente al conjunto de clientes conocidos los fraudes ocurridos hasta el día $t-8$.

El entrenamiento incluyó 4.111 clientes diferentes, de los cuales 1.166 presentaron al menos un fraude. El filtrado redujo las operaciones de validación de 83.437 a 60.655 y los fraudes de 920 a 160. En prueba, las operaciones se redujeron de 85.380 a 58.305 y los fraudes de 920 a 113. En los conjuntos filtrados se identificaron 142 clientes distintos con fraude en validación y 104 en prueba.

La proporción de operaciones fraudulentas resultante fue de 1,280\% en entrenamiento, 0,264\% en validación y 0,194\% en prueba. Estas prevalencias reflejan, entre otros factores, el filtrado aplicado; por lo tanto, no deben compararse como si fueran muestras aleatorias equivalentes.

\subsection{Preprocesamiento para el modelado}
Se excluyeron \texttt{fraud}, \texttt{customer} y \texttt{step} de las variables predictoras: la primera es la etiqueta, mientras que las otras dos se conservaron para organizar el análisis temporal y calcular las métricas por cliente. Se utilizó un \texttt{ColumnTransformer} con \texttt{StandardScaler} para las variables numéricas y \texttt{OneHotEncoder} para \texttt{merchant} y \texttt{category}. El transformador se ajustó exclusivamente con el conjunto de entrenamiento y luego se aplicó a validación y prueba, con tratamiento de categorías no observadas durante el ajuste.
\section{Modelado y evaluación}

\subsection{Modelos considerados}
Se evaluaron cuatro alternativas de clasificación supervisada. La regresión logística se adoptó como \textit{baseline}, por ser una referencia sencilla de interpretar. \textit{Random Forest} y XGBoost permitieron evaluar modelos basados en \textit{árboles}, capaces de representar relaciones no lineales. Por último, se entrenó una red neuronal \textit{Feed-Forward} con dos capas ocultas de 500 neuronas y activación ReLU, \textit{dropout} de 0,2 después de cada capa y salida sigmoide.

La red neuronal se optimizó con Adam y entropía cruzada binaria; se estableció un máximo de 100 \textit{epochs}, lotes de 128 ejemplos y detención temprana con paciencia de cinco \textit{epochs} según el PR-AUC de validación. Los modelos restantes se ejecutaron con los parámetros predefinidos por scikit-learn sin una búsqueda sistemática de hiperparámetros. Cada algoritmo se entrenó en dos variantes: con y sin las dieciocho variables históricas.

\subsection{Métricas y criterio de selección}
La evaluación consideró tres indicadores complementarios. \textit{ROC-AUC} resume la capacidad de ordenar operaciones fraudulentas por encima de las legítimas. \textit{Average Precision} integra el comportamiento de precisión y \textit{recall} a distintos umbrales, y resulta particularmente informativa para un problema de clases desbalanceadas \cite{handbook}. Su expresión es:
\begin{equation}
\label{eq:ap}
AP=\sum_{i}(R_i-R_{i-1})P_i,
\end{equation}
donde $P_i$ y $R_i$ son la precisión y el \textit{recall} de los puntos evaluados.

Para representar la restricción operativa se utilizó CP@$k$. En cada día se agrupan las operaciones por cliente, se conserva el mayor puntaje de riesgo asignado y se seleccionan los $k$ clientes más sospechosos. Un cliente con una o varias operaciones fraudulentas cuenta como una detección correcta. Después de identificarlo, se retira de las evaluaciones de los días posteriores. La métrica general es el promedio de las precisiones diarias:
\begin{equation}
\label{eq:cp}
CP@k=\frac{1}{D}\sum_{d=1}^{D}\frac{C_{d,k}}{k},
\end{equation}
donde $D$ es el número de días y $C_{d,k}$ el número de clientes comprometidos correctamente detectados entre los primeros $k$ del día $d$.

El anteproyecto propuso CP@100 como KPI principal, siguiendo el ejemplo del manual. Sin embargo, en esta simulación el número de clientes comprometidos por día es reducido: en los 23 días de validación hay solamente 142 clientes distintos con fraude. Con una política de retirar clientes ya detectados, el límite de CP@100 es $142/(23\times100)=0,0617$. La regresión logística alcanzó ese valor, lo que hizo poco informativa esta capacidad para distinguir alternativas en la etapa de selección. Por ello se utilizó CP@20 como criterio operativo prioritario, manteniendo AP y ROC-AUC como medidas de apoyo.

\subsection{Resultados en validación}
Los resultados obtenidos con los ocho experimentos se presentan en el Cuadro~\ref{tab:validacion}. Todas las cifras proceden de las salidas calculadas del notebook sobre el conjunto de validación filtrado, correspondiente a los días 127--149.

\begin{table}[!htbp]
\centering
\small
\caption{Comparación de modelos sobre validación}
\label{tab:validacion}
\begin{tabular}{@{}llrrr@{}}
\toprule
\textbf{Modelo} & \textbf{Variables históricas} & \textbf{CP@20} & \textbf{AP} & \textbf{ROC-AUC} \\
\midrule
Regresión logística & Sin & 0,2826 & 0,7555 & 0,9975 \\
Regresión logística & Con & 0,2848 & 0,7985 & 0,9978 \\
Random Forest & Sin & 0,2696 & 0,7091 & 0,9303 \\
Random Forest & Con & 0,2957 & 0,8611 & 0,9959 \\
XGBoost & Sin & 0,2913 & 0,7891 & 0,9984 \\
XGBoost & Con & 0,2870 & 0,7075 & 0,9797 \\
Red neuronal & Sin & 0,2913 & 0,8318 & 0,9984 \\
\textbf{Red neuronal} & \textbf{Con} & \textbf{0,3000} & \textbf{0,8890} & \textbf{0,9993} \\
\bottomrule
\end{tabular}
\par\smallskip
{\footnotesize Valores redondeados a cuatro decimales.}
\end{table}

La ingeniería de variables mejoró la AP en la regresión logística, Random Forest y la red neuronal, aunque no ocurrió lo mismo con XGBoost. En particular, la AP de la red neuronal pasó de 0,8318 sin antecedentes a 0,8890 con ellos, una diferencia de aproximadamente 0,0573 puntos. Su CP@20 aumentó de 0,2913 a 0,3000. En cambio, XGBoost presentó una disminución de AP y CP@20 al incorporar las nuevas variables. Por lo tanto, los resultados respaldan la utilidad de las características históricas para algunos algoritmos, pero no una mejora generalizada.

% \begin{figure}[!htbp]
% \centering
% \begin{tikzpicture}
% \begin{axis}[
%     width=0.91\linewidth,height=6.0cm,
%     ybar,bar width=9pt,
%     ymin=0,ymax=1,
%     ylabel={Average Precision},
%     symbolic x coords={Logística,Forest,XGBoost,Red neuronal},
%     xtick=data,
%     x tick label style={font=\small},
%     y tick label style={font=\small},
%     legend style={at={(0.5,1.14)},anchor=south,legend columns=2,draw=none,font=\small},
%     grid=major,grid style={gray!15},
%     axis x line*=bottom,axis y line*=left,
%     enlarge x limits=0.15
% ]
% \addplot[fill=blue!50,draw=blue!80!black] coordinates {
%  (Logística,0.755506) (Forest,0.709099) (XGBoost,0.789103) (Red neuronal,0.831752)};
% \addplot[fill=pink!50,draw=pink!80!black] coordinates {
%  (Logística,0.798490) (Forest,0.861130) (XGBoost,0.707471) (Red neuronal,0.889045)};
% \legend{Sin variables históricas,Con variables históricas}
% \end{axis}
% \end{tikzpicture}
% \caption{Average Precision con y sin ingeniería de variables, calculada en validación.}
% \label{fig:ap}
% \end{figure}

\begin{figure}[!htbp]
    \centering
    \resizebox{1\textwidth}{!}{%% Creator: Matplotlib, PGF backend
%%
%% To include the figure in your LaTeX document, write
%%   \input{<filename>.pgf}
%%
%% Make sure the required packages are loaded in your preamble
%%   \usepackage{pgf}
%%
%% Also ensure that all the required font packages are loaded; for instance,
%% the lmodern package is sometimes necessary when using math font.
%%   \usepackage{lmodern}
%%
%% Figures using additional raster images can only be included by \input if
%% they are in the same directory as the main LaTeX file. For loading figures
%% from other directories you can use the `import` package
%%   \usepackage{import}
%%
%% and then include the figures with
%%   \import{<path to file>}{<filename>.pgf}
%%
%% Matplotlib used the following preamble
%%   \def\mathdefault#1{#1}
%%   \everymath=\expandafter{\the\everymath\displaystyle}
%%   \IfFileExists{scrextend.sty}{
%%     \usepackage[fontsize=10.000000pt]{scrextend}
%%   }{
%%     \renewcommand{\normalsize}{\fontsize{10.000000}{12.000000}\selectfont}
%%     \normalsize
%%   }
%%   
%%   \ifdefined\pdftexversion\else  % non-pdftex case.
%%     \usepackage{fontspec}
%%     \setmainfont{DejaVuSerif.ttf}[Path=\detokenize{/home/ginoanci/Documentos/Diplomaturas/Trabajos_Practicos/Practicos_Diplo_CD/TF/.venv/lib/python3.13/site-packages/matplotlib/mpl-data/fonts/ttf/}]
%%     \setsansfont{DejaVuSans.ttf}[Path=\detokenize{/home/ginoanci/Documentos/Diplomaturas/Trabajos_Practicos/Practicos_Diplo_CD/TF/.venv/lib/python3.13/site-packages/matplotlib/mpl-data/fonts/ttf/}]
%%     \setmonofont{DejaVuSansMono.ttf}[Path=\detokenize{/home/ginoanci/Documentos/Diplomaturas/Trabajos_Practicos/Practicos_Diplo_CD/TF/.venv/lib/python3.13/site-packages/matplotlib/mpl-data/fonts/ttf/}]
%%   \fi
%%   \makeatletter\@ifpackageloaded{underscore}{}{\usepackage[strings]{underscore}}\makeatother
%%
\begingroup%
\makeatletter%
\begin{pgfpicture}%
\pgfpathrectangle{\pgfpointorigin}{\pgfqpoint{7.900879in}{2.678502in}}%
\pgfusepath{use as bounding box, clip}%
\begin{pgfscope}%
\pgfsetbuttcap%
\pgfsetmiterjoin%
\definecolor{currentfill}{rgb}{1.000000,1.000000,1.000000}%
\pgfsetfillcolor{currentfill}%
\pgfsetlinewidth{0.000000pt}%
\definecolor{currentstroke}{rgb}{1.000000,1.000000,1.000000}%
\pgfsetstrokecolor{currentstroke}%
\pgfsetdash{}{0pt}%
\pgfpathmoveto{\pgfqpoint{0.000000in}{0.000000in}}%
\pgfpathlineto{\pgfqpoint{7.900879in}{0.000000in}}%
\pgfpathlineto{\pgfqpoint{7.900879in}{2.678502in}}%
\pgfpathlineto{\pgfqpoint{0.000000in}{2.678502in}}%
\pgfpathlineto{\pgfqpoint{0.000000in}{0.000000in}}%
\pgfpathclose%
\pgfusepath{fill}%
\end{pgfscope}%
\begin{pgfscope}%
\pgfsetbuttcap%
\pgfsetmiterjoin%
\definecolor{currentfill}{rgb}{1.000000,1.000000,1.000000}%
\pgfsetfillcolor{currentfill}%
\pgfsetlinewidth{0.000000pt}%
\definecolor{currentstroke}{rgb}{0.000000,0.000000,0.000000}%
\pgfsetstrokecolor{currentstroke}%
\pgfsetstrokeopacity{0.000000}%
\pgfsetdash{}{0pt}%
\pgfpathmoveto{\pgfqpoint{0.626435in}{0.341672in}}%
\pgfpathlineto{\pgfqpoint{7.800879in}{0.341672in}}%
\pgfpathlineto{\pgfqpoint{7.800879in}{2.347128in}}%
\pgfpathlineto{\pgfqpoint{0.626435in}{2.347128in}}%
\pgfpathlineto{\pgfqpoint{0.626435in}{0.341672in}}%
\pgfpathclose%
\pgfusepath{fill}%
\end{pgfscope}%
\begin{pgfscope}%
\pgfpathrectangle{\pgfqpoint{0.626435in}{0.341672in}}{\pgfqpoint{7.174444in}{2.005456in}}%
\pgfusepath{clip}%
\pgfsetbuttcap%
\pgfsetmiterjoin%
\definecolor{currentfill}{rgb}{0.560784,0.658824,0.749020}%
\pgfsetfillcolor{currentfill}%
\pgfsetlinewidth{0.803000pt}%
\definecolor{currentstroke}{rgb}{0.435294,0.529412,0.611765}%
\pgfsetstrokecolor{currentstroke}%
\pgfsetdash{}{0pt}%
\pgfpathmoveto{\pgfqpoint{0.952546in}{0.341672in}}%
\pgfpathlineto{\pgfqpoint{1.301951in}{0.341672in}}%
\pgfpathlineto{\pgfqpoint{1.301951in}{1.856806in}}%
\pgfpathlineto{\pgfqpoint{0.952546in}{1.856806in}}%
\pgfpathlineto{\pgfqpoint{0.952546in}{0.341672in}}%
\pgfpathclose%
\pgfusepath{stroke,fill}%
\end{pgfscope}%
\begin{pgfscope}%
\pgfpathrectangle{\pgfqpoint{0.626435in}{0.341672in}}{\pgfqpoint{7.174444in}{2.005456in}}%
\pgfusepath{clip}%
\pgfsetbuttcap%
\pgfsetmiterjoin%
\definecolor{currentfill}{rgb}{0.560784,0.658824,0.749020}%
\pgfsetfillcolor{currentfill}%
\pgfsetlinewidth{0.803000pt}%
\definecolor{currentstroke}{rgb}{0.435294,0.529412,0.611765}%
\pgfsetstrokecolor{currentstroke}%
\pgfsetdash{}{0pt}%
\pgfpathmoveto{\pgfqpoint{2.893684in}{0.341672in}}%
\pgfpathlineto{\pgfqpoint{3.243088in}{0.341672in}}%
\pgfpathlineto{\pgfqpoint{3.243088in}{1.763739in}}%
\pgfpathlineto{\pgfqpoint{2.893684in}{1.763739in}}%
\pgfpathlineto{\pgfqpoint{2.893684in}{0.341672in}}%
\pgfpathclose%
\pgfusepath{stroke,fill}%
\end{pgfscope}%
\begin{pgfscope}%
\pgfpathrectangle{\pgfqpoint{0.626435in}{0.341672in}}{\pgfqpoint{7.174444in}{2.005456in}}%
\pgfusepath{clip}%
\pgfsetbuttcap%
\pgfsetmiterjoin%
\definecolor{currentfill}{rgb}{0.560784,0.658824,0.749020}%
\pgfsetfillcolor{currentfill}%
\pgfsetlinewidth{0.803000pt}%
\definecolor{currentstroke}{rgb}{0.435294,0.529412,0.611765}%
\pgfsetstrokecolor{currentstroke}%
\pgfsetdash{}{0pt}%
\pgfpathmoveto{\pgfqpoint{4.834821in}{0.341672in}}%
\pgfpathlineto{\pgfqpoint{5.184226in}{0.341672in}}%
\pgfpathlineto{\pgfqpoint{5.184226in}{1.924184in}}%
\pgfpathlineto{\pgfqpoint{4.834821in}{1.924184in}}%
\pgfpathlineto{\pgfqpoint{4.834821in}{0.341672in}}%
\pgfpathclose%
\pgfusepath{stroke,fill}%
\end{pgfscope}%
\begin{pgfscope}%
\pgfpathrectangle{\pgfqpoint{0.626435in}{0.341672in}}{\pgfqpoint{7.174444in}{2.005456in}}%
\pgfusepath{clip}%
\pgfsetbuttcap%
\pgfsetmiterjoin%
\definecolor{currentfill}{rgb}{0.560784,0.658824,0.749020}%
\pgfsetfillcolor{currentfill}%
\pgfsetlinewidth{0.803000pt}%
\definecolor{currentstroke}{rgb}{0.435294,0.529412,0.611765}%
\pgfsetstrokecolor{currentstroke}%
\pgfsetdash{}{0pt}%
\pgfpathmoveto{\pgfqpoint{6.775959in}{0.341672in}}%
\pgfpathlineto{\pgfqpoint{7.125364in}{0.341672in}}%
\pgfpathlineto{\pgfqpoint{7.125364in}{2.008116in}}%
\pgfpathlineto{\pgfqpoint{6.775959in}{2.008116in}}%
\pgfpathlineto{\pgfqpoint{6.775959in}{0.341672in}}%
\pgfpathclose%
\pgfusepath{stroke,fill}%
\end{pgfscope}%
\begin{pgfscope}%
\pgfpathrectangle{\pgfqpoint{0.626435in}{0.341672in}}{\pgfqpoint{7.174444in}{2.005456in}}%
\pgfusepath{clip}%
\pgfsetbuttcap%
\pgfsetmiterjoin%
\definecolor{currentfill}{rgb}{0.850980,0.549020,0.584314}%
\pgfsetfillcolor{currentfill}%
\pgfsetlinewidth{0.803000pt}%
\definecolor{currentstroke}{rgb}{0.709804,0.431373,0.470588}%
\pgfsetstrokecolor{currentstroke}%
\pgfsetdash{}{0pt}%
\pgfpathmoveto{\pgfqpoint{1.301951in}{0.341672in}}%
\pgfpathlineto{\pgfqpoint{1.651356in}{0.341672in}}%
\pgfpathlineto{\pgfqpoint{1.651356in}{1.944032in}}%
\pgfpathlineto{\pgfqpoint{1.301951in}{1.944032in}}%
\pgfpathlineto{\pgfqpoint{1.301951in}{0.341672in}}%
\pgfpathclose%
\pgfusepath{stroke,fill}%
\end{pgfscope}%
\begin{pgfscope}%
\pgfpathrectangle{\pgfqpoint{0.626435in}{0.341672in}}{\pgfqpoint{7.174444in}{2.005456in}}%
\pgfusepath{clip}%
\pgfsetbuttcap%
\pgfsetmiterjoin%
\definecolor{currentfill}{rgb}{0.850980,0.549020,0.584314}%
\pgfsetfillcolor{currentfill}%
\pgfsetlinewidth{0.803000pt}%
\definecolor{currentstroke}{rgb}{0.709804,0.431373,0.470588}%
\pgfsetstrokecolor{currentstroke}%
\pgfsetdash{}{0pt}%
\pgfpathmoveto{\pgfqpoint{3.243088in}{0.341672in}}%
\pgfpathlineto{\pgfqpoint{3.592493in}{0.341672in}}%
\pgfpathlineto{\pgfqpoint{3.592493in}{2.068631in}}%
\pgfpathlineto{\pgfqpoint{3.243088in}{2.068631in}}%
\pgfpathlineto{\pgfqpoint{3.243088in}{0.341672in}}%
\pgfpathclose%
\pgfusepath{stroke,fill}%
\end{pgfscope}%
\begin{pgfscope}%
\pgfpathrectangle{\pgfqpoint{0.626435in}{0.341672in}}{\pgfqpoint{7.174444in}{2.005456in}}%
\pgfusepath{clip}%
\pgfsetbuttcap%
\pgfsetmiterjoin%
\definecolor{currentfill}{rgb}{0.850980,0.549020,0.584314}%
\pgfsetfillcolor{currentfill}%
\pgfsetlinewidth{0.803000pt}%
\definecolor{currentstroke}{rgb}{0.709804,0.431373,0.470588}%
\pgfsetstrokecolor{currentstroke}%
\pgfsetdash{}{0pt}%
\pgfpathmoveto{\pgfqpoint{5.184226in}{0.341672in}}%
\pgfpathlineto{\pgfqpoint{5.533631in}{0.341672in}}%
\pgfpathlineto{\pgfqpoint{5.533631in}{1.760474in}}%
\pgfpathlineto{\pgfqpoint{5.184226in}{1.760474in}}%
\pgfpathlineto{\pgfqpoint{5.184226in}{0.341672in}}%
\pgfpathclose%
\pgfusepath{stroke,fill}%
\end{pgfscope}%
\begin{pgfscope}%
\pgfpathrectangle{\pgfqpoint{0.626435in}{0.341672in}}{\pgfqpoint{7.174444in}{2.005456in}}%
\pgfusepath{clip}%
\pgfsetbuttcap%
\pgfsetmiterjoin%
\definecolor{currentfill}{rgb}{0.850980,0.549020,0.584314}%
\pgfsetfillcolor{currentfill}%
\pgfsetlinewidth{0.803000pt}%
\definecolor{currentstroke}{rgb}{0.709804,0.431373,0.470588}%
\pgfsetstrokecolor{currentstroke}%
\pgfsetdash{}{0pt}%
\pgfpathmoveto{\pgfqpoint{7.125364in}{0.341672in}}%
\pgfpathlineto{\pgfqpoint{7.474768in}{0.341672in}}%
\pgfpathlineto{\pgfqpoint{7.474768in}{2.112941in}}%
\pgfpathlineto{\pgfqpoint{7.125364in}{2.112941in}}%
\pgfpathlineto{\pgfqpoint{7.125364in}{0.341672in}}%
\pgfpathclose%
\pgfusepath{stroke,fill}%
\end{pgfscope}%
\begin{pgfscope}%
\pgfsetbuttcap%
\pgfsetroundjoin%
\definecolor{currentfill}{rgb}{0.000000,0.000000,0.000000}%
\pgfsetfillcolor{currentfill}%
\pgfsetlinewidth{0.803000pt}%
\definecolor{currentstroke}{rgb}{0.000000,0.000000,0.000000}%
\pgfsetstrokecolor{currentstroke}%
\pgfsetdash{}{0pt}%
\pgfsys@defobject{currentmarker}{\pgfqpoint{0.000000in}{-0.048611in}}{\pgfqpoint{0.000000in}{0.000000in}}{%
\pgfpathmoveto{\pgfqpoint{0.000000in}{0.000000in}}%
\pgfpathlineto{\pgfqpoint{0.000000in}{-0.048611in}}%
\pgfusepath{stroke,fill}%
}%
\begin{pgfscope}%
\pgfsys@transformshift{1.301951in}{0.341672in}%
\pgfsys@useobject{currentmarker}{}%
\end{pgfscope}%
\end{pgfscope}%
\begin{pgfscope}%
\definecolor{textcolor}{rgb}{0.000000,0.000000,0.000000}%
\pgfsetstrokecolor{textcolor}%
\pgfsetfillcolor{textcolor}%
\pgftext[x=1.301951in,y=0.244450in,,top]{\color{textcolor}{\sffamily\fontsize{10.000000}{12.000000}\selectfont\catcode`\^=\active\def^{\ifmmode\sp\else\^{}\fi}\catcode`\%=\active\def%{\%}Logística}}%
\end{pgfscope}%
\begin{pgfscope}%
\pgfsetbuttcap%
\pgfsetroundjoin%
\definecolor{currentfill}{rgb}{0.000000,0.000000,0.000000}%
\pgfsetfillcolor{currentfill}%
\pgfsetlinewidth{0.803000pt}%
\definecolor{currentstroke}{rgb}{0.000000,0.000000,0.000000}%
\pgfsetstrokecolor{currentstroke}%
\pgfsetdash{}{0pt}%
\pgfsys@defobject{currentmarker}{\pgfqpoint{0.000000in}{-0.048611in}}{\pgfqpoint{0.000000in}{0.000000in}}{%
\pgfpathmoveto{\pgfqpoint{0.000000in}{0.000000in}}%
\pgfpathlineto{\pgfqpoint{0.000000in}{-0.048611in}}%
\pgfusepath{stroke,fill}%
}%
\begin{pgfscope}%
\pgfsys@transformshift{3.243088in}{0.341672in}%
\pgfsys@useobject{currentmarker}{}%
\end{pgfscope}%
\end{pgfscope}%
\begin{pgfscope}%
\definecolor{textcolor}{rgb}{0.000000,0.000000,0.000000}%
\pgfsetstrokecolor{textcolor}%
\pgfsetfillcolor{textcolor}%
\pgftext[x=3.243088in,y=0.244450in,,top]{\color{textcolor}{\sffamily\fontsize{10.000000}{12.000000}\selectfont\catcode`\^=\active\def^{\ifmmode\sp\else\^{}\fi}\catcode`\%=\active\def%{\%}Random Forest}}%
\end{pgfscope}%
\begin{pgfscope}%
\pgfsetbuttcap%
\pgfsetroundjoin%
\definecolor{currentfill}{rgb}{0.000000,0.000000,0.000000}%
\pgfsetfillcolor{currentfill}%
\pgfsetlinewidth{0.803000pt}%
\definecolor{currentstroke}{rgb}{0.000000,0.000000,0.000000}%
\pgfsetstrokecolor{currentstroke}%
\pgfsetdash{}{0pt}%
\pgfsys@defobject{currentmarker}{\pgfqpoint{0.000000in}{-0.048611in}}{\pgfqpoint{0.000000in}{0.000000in}}{%
\pgfpathmoveto{\pgfqpoint{0.000000in}{0.000000in}}%
\pgfpathlineto{\pgfqpoint{0.000000in}{-0.048611in}}%
\pgfusepath{stroke,fill}%
}%
\begin{pgfscope}%
\pgfsys@transformshift{5.184226in}{0.341672in}%
\pgfsys@useobject{currentmarker}{}%
\end{pgfscope}%
\end{pgfscope}%
\begin{pgfscope}%
\definecolor{textcolor}{rgb}{0.000000,0.000000,0.000000}%
\pgfsetstrokecolor{textcolor}%
\pgfsetfillcolor{textcolor}%
\pgftext[x=5.184226in,y=0.244450in,,top]{\color{textcolor}{\sffamily\fontsize{10.000000}{12.000000}\selectfont\catcode`\^=\active\def^{\ifmmode\sp\else\^{}\fi}\catcode`\%=\active\def%{\%}XGBoost}}%
\end{pgfscope}%
\begin{pgfscope}%
\pgfsetbuttcap%
\pgfsetroundjoin%
\definecolor{currentfill}{rgb}{0.000000,0.000000,0.000000}%
\pgfsetfillcolor{currentfill}%
\pgfsetlinewidth{0.803000pt}%
\definecolor{currentstroke}{rgb}{0.000000,0.000000,0.000000}%
\pgfsetstrokecolor{currentstroke}%
\pgfsetdash{}{0pt}%
\pgfsys@defobject{currentmarker}{\pgfqpoint{0.000000in}{-0.048611in}}{\pgfqpoint{0.000000in}{0.000000in}}{%
\pgfpathmoveto{\pgfqpoint{0.000000in}{0.000000in}}%
\pgfpathlineto{\pgfqpoint{0.000000in}{-0.048611in}}%
\pgfusepath{stroke,fill}%
}%
\begin{pgfscope}%
\pgfsys@transformshift{7.125364in}{0.341672in}%
\pgfsys@useobject{currentmarker}{}%
\end{pgfscope}%
\end{pgfscope}%
\begin{pgfscope}%
\definecolor{textcolor}{rgb}{0.000000,0.000000,0.000000}%
\pgfsetstrokecolor{textcolor}%
\pgfsetfillcolor{textcolor}%
\pgftext[x=7.125364in,y=0.244450in,,top]{\color{textcolor}{\sffamily\fontsize{10.000000}{12.000000}\selectfont\catcode`\^=\active\def^{\ifmmode\sp\else\^{}\fi}\catcode`\%=\active\def%{\%}Red neuronal}}%
\end{pgfscope}%
\begin{pgfscope}%
\pgfpathrectangle{\pgfqpoint{0.626435in}{0.341672in}}{\pgfqpoint{7.174444in}{2.005456in}}%
\pgfusepath{clip}%
\pgfsetrectcap%
\pgfsetroundjoin%
\pgfsetlinewidth{0.803000pt}%
\definecolor{currentstroke}{rgb}{0.690196,0.690196,0.690196}%
\pgfsetstrokecolor{currentstroke}%
\pgfsetstrokeopacity{0.200000}%
\pgfsetdash{}{0pt}%
\pgfpathmoveto{\pgfqpoint{0.626435in}{0.341672in}}%
\pgfpathlineto{\pgfqpoint{7.800879in}{0.341672in}}%
\pgfusepath{stroke}%
\end{pgfscope}%
\begin{pgfscope}%
\pgfsetbuttcap%
\pgfsetroundjoin%
\definecolor{currentfill}{rgb}{0.000000,0.000000,0.000000}%
\pgfsetfillcolor{currentfill}%
\pgfsetlinewidth{0.803000pt}%
\definecolor{currentstroke}{rgb}{0.000000,0.000000,0.000000}%
\pgfsetstrokecolor{currentstroke}%
\pgfsetdash{}{0pt}%
\pgfsys@defobject{currentmarker}{\pgfqpoint{-0.048611in}{0.000000in}}{\pgfqpoint{-0.000000in}{0.000000in}}{%
\pgfpathmoveto{\pgfqpoint{-0.000000in}{0.000000in}}%
\pgfpathlineto{\pgfqpoint{-0.048611in}{0.000000in}}%
\pgfusepath{stroke,fill}%
}%
\begin{pgfscope}%
\pgfsys@transformshift{0.626435in}{0.341672in}%
\pgfsys@useobject{currentmarker}{}%
\end{pgfscope}%
\end{pgfscope}%
\begin{pgfscope}%
\definecolor{textcolor}{rgb}{0.000000,0.000000,0.000000}%
\pgfsetstrokecolor{textcolor}%
\pgfsetfillcolor{textcolor}%
\pgftext[x=0.308333in, y=0.288911in, left, base]{\color{textcolor}{\sffamily\fontsize{10.000000}{12.000000}\selectfont\catcode`\^=\active\def^{\ifmmode\sp\else\^{}\fi}\catcode`\%=\active\def%{\%}0.0}}%
\end{pgfscope}%
\begin{pgfscope}%
\pgfpathrectangle{\pgfqpoint{0.626435in}{0.341672in}}{\pgfqpoint{7.174444in}{2.005456in}}%
\pgfusepath{clip}%
\pgfsetrectcap%
\pgfsetroundjoin%
\pgfsetlinewidth{0.803000pt}%
\definecolor{currentstroke}{rgb}{0.690196,0.690196,0.690196}%
\pgfsetstrokecolor{currentstroke}%
\pgfsetstrokeopacity{0.200000}%
\pgfsetdash{}{0pt}%
\pgfpathmoveto{\pgfqpoint{0.626435in}{0.742763in}}%
\pgfpathlineto{\pgfqpoint{7.800879in}{0.742763in}}%
\pgfusepath{stroke}%
\end{pgfscope}%
\begin{pgfscope}%
\pgfsetbuttcap%
\pgfsetroundjoin%
\definecolor{currentfill}{rgb}{0.000000,0.000000,0.000000}%
\pgfsetfillcolor{currentfill}%
\pgfsetlinewidth{0.803000pt}%
\definecolor{currentstroke}{rgb}{0.000000,0.000000,0.000000}%
\pgfsetstrokecolor{currentstroke}%
\pgfsetdash{}{0pt}%
\pgfsys@defobject{currentmarker}{\pgfqpoint{-0.048611in}{0.000000in}}{\pgfqpoint{-0.000000in}{0.000000in}}{%
\pgfpathmoveto{\pgfqpoint{-0.000000in}{0.000000in}}%
\pgfpathlineto{\pgfqpoint{-0.048611in}{0.000000in}}%
\pgfusepath{stroke,fill}%
}%
\begin{pgfscope}%
\pgfsys@transformshift{0.626435in}{0.742763in}%
\pgfsys@useobject{currentmarker}{}%
\end{pgfscope}%
\end{pgfscope}%
\begin{pgfscope}%
\definecolor{textcolor}{rgb}{0.000000,0.000000,0.000000}%
\pgfsetstrokecolor{textcolor}%
\pgfsetfillcolor{textcolor}%
\pgftext[x=0.308333in, y=0.690002in, left, base]{\color{textcolor}{\sffamily\fontsize{10.000000}{12.000000}\selectfont\catcode`\^=\active\def^{\ifmmode\sp\else\^{}\fi}\catcode`\%=\active\def%{\%}0.2}}%
\end{pgfscope}%
\begin{pgfscope}%
\pgfpathrectangle{\pgfqpoint{0.626435in}{0.341672in}}{\pgfqpoint{7.174444in}{2.005456in}}%
\pgfusepath{clip}%
\pgfsetrectcap%
\pgfsetroundjoin%
\pgfsetlinewidth{0.803000pt}%
\definecolor{currentstroke}{rgb}{0.690196,0.690196,0.690196}%
\pgfsetstrokecolor{currentstroke}%
\pgfsetstrokeopacity{0.200000}%
\pgfsetdash{}{0pt}%
\pgfpathmoveto{\pgfqpoint{0.626435in}{1.143855in}}%
\pgfpathlineto{\pgfqpoint{7.800879in}{1.143855in}}%
\pgfusepath{stroke}%
\end{pgfscope}%
\begin{pgfscope}%
\pgfsetbuttcap%
\pgfsetroundjoin%
\definecolor{currentfill}{rgb}{0.000000,0.000000,0.000000}%
\pgfsetfillcolor{currentfill}%
\pgfsetlinewidth{0.803000pt}%
\definecolor{currentstroke}{rgb}{0.000000,0.000000,0.000000}%
\pgfsetstrokecolor{currentstroke}%
\pgfsetdash{}{0pt}%
\pgfsys@defobject{currentmarker}{\pgfqpoint{-0.048611in}{0.000000in}}{\pgfqpoint{-0.000000in}{0.000000in}}{%
\pgfpathmoveto{\pgfqpoint{-0.000000in}{0.000000in}}%
\pgfpathlineto{\pgfqpoint{-0.048611in}{0.000000in}}%
\pgfusepath{stroke,fill}%
}%
\begin{pgfscope}%
\pgfsys@transformshift{0.626435in}{1.143855in}%
\pgfsys@useobject{currentmarker}{}%
\end{pgfscope}%
\end{pgfscope}%
\begin{pgfscope}%
\definecolor{textcolor}{rgb}{0.000000,0.000000,0.000000}%
\pgfsetstrokecolor{textcolor}%
\pgfsetfillcolor{textcolor}%
\pgftext[x=0.308333in, y=1.091093in, left, base]{\color{textcolor}{\sffamily\fontsize{10.000000}{12.000000}\selectfont\catcode`\^=\active\def^{\ifmmode\sp\else\^{}\fi}\catcode`\%=\active\def%{\%}0.4}}%
\end{pgfscope}%
\begin{pgfscope}%
\pgfpathrectangle{\pgfqpoint{0.626435in}{0.341672in}}{\pgfqpoint{7.174444in}{2.005456in}}%
\pgfusepath{clip}%
\pgfsetrectcap%
\pgfsetroundjoin%
\pgfsetlinewidth{0.803000pt}%
\definecolor{currentstroke}{rgb}{0.690196,0.690196,0.690196}%
\pgfsetstrokecolor{currentstroke}%
\pgfsetstrokeopacity{0.200000}%
\pgfsetdash{}{0pt}%
\pgfpathmoveto{\pgfqpoint{0.626435in}{1.544946in}}%
\pgfpathlineto{\pgfqpoint{7.800879in}{1.544946in}}%
\pgfusepath{stroke}%
\end{pgfscope}%
\begin{pgfscope}%
\pgfsetbuttcap%
\pgfsetroundjoin%
\definecolor{currentfill}{rgb}{0.000000,0.000000,0.000000}%
\pgfsetfillcolor{currentfill}%
\pgfsetlinewidth{0.803000pt}%
\definecolor{currentstroke}{rgb}{0.000000,0.000000,0.000000}%
\pgfsetstrokecolor{currentstroke}%
\pgfsetdash{}{0pt}%
\pgfsys@defobject{currentmarker}{\pgfqpoint{-0.048611in}{0.000000in}}{\pgfqpoint{-0.000000in}{0.000000in}}{%
\pgfpathmoveto{\pgfqpoint{-0.000000in}{0.000000in}}%
\pgfpathlineto{\pgfqpoint{-0.048611in}{0.000000in}}%
\pgfusepath{stroke,fill}%
}%
\begin{pgfscope}%
\pgfsys@transformshift{0.626435in}{1.544946in}%
\pgfsys@useobject{currentmarker}{}%
\end{pgfscope}%
\end{pgfscope}%
\begin{pgfscope}%
\definecolor{textcolor}{rgb}{0.000000,0.000000,0.000000}%
\pgfsetstrokecolor{textcolor}%
\pgfsetfillcolor{textcolor}%
\pgftext[x=0.308333in, y=1.492184in, left, base]{\color{textcolor}{\sffamily\fontsize{10.000000}{12.000000}\selectfont\catcode`\^=\active\def^{\ifmmode\sp\else\^{}\fi}\catcode`\%=\active\def%{\%}0.6}}%
\end{pgfscope}%
\begin{pgfscope}%
\pgfpathrectangle{\pgfqpoint{0.626435in}{0.341672in}}{\pgfqpoint{7.174444in}{2.005456in}}%
\pgfusepath{clip}%
\pgfsetrectcap%
\pgfsetroundjoin%
\pgfsetlinewidth{0.803000pt}%
\definecolor{currentstroke}{rgb}{0.690196,0.690196,0.690196}%
\pgfsetstrokecolor{currentstroke}%
\pgfsetstrokeopacity{0.200000}%
\pgfsetdash{}{0pt}%
\pgfpathmoveto{\pgfqpoint{0.626435in}{1.946037in}}%
\pgfpathlineto{\pgfqpoint{7.800879in}{1.946037in}}%
\pgfusepath{stroke}%
\end{pgfscope}%
\begin{pgfscope}%
\pgfsetbuttcap%
\pgfsetroundjoin%
\definecolor{currentfill}{rgb}{0.000000,0.000000,0.000000}%
\pgfsetfillcolor{currentfill}%
\pgfsetlinewidth{0.803000pt}%
\definecolor{currentstroke}{rgb}{0.000000,0.000000,0.000000}%
\pgfsetstrokecolor{currentstroke}%
\pgfsetdash{}{0pt}%
\pgfsys@defobject{currentmarker}{\pgfqpoint{-0.048611in}{0.000000in}}{\pgfqpoint{-0.000000in}{0.000000in}}{%
\pgfpathmoveto{\pgfqpoint{-0.000000in}{0.000000in}}%
\pgfpathlineto{\pgfqpoint{-0.048611in}{0.000000in}}%
\pgfusepath{stroke,fill}%
}%
\begin{pgfscope}%
\pgfsys@transformshift{0.626435in}{1.946037in}%
\pgfsys@useobject{currentmarker}{}%
\end{pgfscope}%
\end{pgfscope}%
\begin{pgfscope}%
\definecolor{textcolor}{rgb}{0.000000,0.000000,0.000000}%
\pgfsetstrokecolor{textcolor}%
\pgfsetfillcolor{textcolor}%
\pgftext[x=0.308333in, y=1.893276in, left, base]{\color{textcolor}{\sffamily\fontsize{10.000000}{12.000000}\selectfont\catcode`\^=\active\def^{\ifmmode\sp\else\^{}\fi}\catcode`\%=\active\def%{\%}0.8}}%
\end{pgfscope}%
\begin{pgfscope}%
\pgfpathrectangle{\pgfqpoint{0.626435in}{0.341672in}}{\pgfqpoint{7.174444in}{2.005456in}}%
\pgfusepath{clip}%
\pgfsetrectcap%
\pgfsetroundjoin%
\pgfsetlinewidth{0.803000pt}%
\definecolor{currentstroke}{rgb}{0.690196,0.690196,0.690196}%
\pgfsetstrokecolor{currentstroke}%
\pgfsetstrokeopacity{0.200000}%
\pgfsetdash{}{0pt}%
\pgfpathmoveto{\pgfqpoint{0.626435in}{2.347128in}}%
\pgfpathlineto{\pgfqpoint{7.800879in}{2.347128in}}%
\pgfusepath{stroke}%
\end{pgfscope}%
\begin{pgfscope}%
\pgfsetbuttcap%
\pgfsetroundjoin%
\definecolor{currentfill}{rgb}{0.000000,0.000000,0.000000}%
\pgfsetfillcolor{currentfill}%
\pgfsetlinewidth{0.803000pt}%
\definecolor{currentstroke}{rgb}{0.000000,0.000000,0.000000}%
\pgfsetstrokecolor{currentstroke}%
\pgfsetdash{}{0pt}%
\pgfsys@defobject{currentmarker}{\pgfqpoint{-0.048611in}{0.000000in}}{\pgfqpoint{-0.000000in}{0.000000in}}{%
\pgfpathmoveto{\pgfqpoint{-0.000000in}{0.000000in}}%
\pgfpathlineto{\pgfqpoint{-0.048611in}{0.000000in}}%
\pgfusepath{stroke,fill}%
}%
\begin{pgfscope}%
\pgfsys@transformshift{0.626435in}{2.347128in}%
\pgfsys@useobject{currentmarker}{}%
\end{pgfscope}%
\end{pgfscope}%
\begin{pgfscope}%
\definecolor{textcolor}{rgb}{0.000000,0.000000,0.000000}%
\pgfsetstrokecolor{textcolor}%
\pgfsetfillcolor{textcolor}%
\pgftext[x=0.308333in, y=2.294367in, left, base]{\color{textcolor}{\sffamily\fontsize{10.000000}{12.000000}\selectfont\catcode`\^=\active\def^{\ifmmode\sp\else\^{}\fi}\catcode`\%=\active\def%{\%}1.0}}%
\end{pgfscope}%
\begin{pgfscope}%
\definecolor{textcolor}{rgb}{0.000000,0.000000,0.000000}%
\pgfsetstrokecolor{textcolor}%
\pgfsetfillcolor{textcolor}%
\pgftext[x=0.252778in,y=1.344400in,,bottom,rotate=90.000000]{\color{textcolor}{\sffamily\fontsize{11.000000}{13.200000}\selectfont\catcode`\^=\active\def^{\ifmmode\sp\else\^{}\fi}\catcode`\%=\active\def%{\%}Average Precision}}%
\end{pgfscope}%
\begin{pgfscope}%
\pgfsetrectcap%
\pgfsetmiterjoin%
\pgfsetlinewidth{0.803000pt}%
\definecolor{currentstroke}{rgb}{0.000000,0.000000,0.000000}%
\pgfsetstrokecolor{currentstroke}%
\pgfsetdash{}{0pt}%
\pgfpathmoveto{\pgfqpoint{0.626435in}{0.341672in}}%
\pgfpathlineto{\pgfqpoint{0.626435in}{2.347128in}}%
\pgfusepath{stroke}%
\end{pgfscope}%
\begin{pgfscope}%
\pgfsetrectcap%
\pgfsetmiterjoin%
\pgfsetlinewidth{0.803000pt}%
\definecolor{currentstroke}{rgb}{0.000000,0.000000,0.000000}%
\pgfsetstrokecolor{currentstroke}%
\pgfsetdash{}{0pt}%
\pgfpathmoveto{\pgfqpoint{7.800879in}{0.341672in}}%
\pgfpathlineto{\pgfqpoint{7.800879in}{2.347128in}}%
\pgfusepath{stroke}%
\end{pgfscope}%
\begin{pgfscope}%
\pgfsetrectcap%
\pgfsetmiterjoin%
\pgfsetlinewidth{0.803000pt}%
\definecolor{currentstroke}{rgb}{0.000000,0.000000,0.000000}%
\pgfsetstrokecolor{currentstroke}%
\pgfsetdash{}{0pt}%
\pgfpathmoveto{\pgfqpoint{0.626435in}{0.341672in}}%
\pgfpathlineto{\pgfqpoint{7.800879in}{0.341672in}}%
\pgfusepath{stroke}%
\end{pgfscope}%
\begin{pgfscope}%
\pgfsetrectcap%
\pgfsetmiterjoin%
\pgfsetlinewidth{0.803000pt}%
\definecolor{currentstroke}{rgb}{0.000000,0.000000,0.000000}%
\pgfsetstrokecolor{currentstroke}%
\pgfsetdash{}{0pt}%
\pgfpathmoveto{\pgfqpoint{0.626435in}{2.347128in}}%
\pgfpathlineto{\pgfqpoint{7.800879in}{2.347128in}}%
\pgfusepath{stroke}%
\end{pgfscope}%
\begin{pgfscope}%
\definecolor{textcolor}{rgb}{0.000000,0.000000,0.000000}%
\pgfsetstrokecolor{textcolor}%
\pgfsetfillcolor{textcolor}%
\pgftext[x=1.127248in,y=1.898473in,,bottom]{\color{textcolor}{\sffamily\fontsize{9.000000}{10.800000}\selectfont\catcode`\^=\active\def^{\ifmmode\sp\else\^{}\fi}\catcode`\%=\active\def%{\%}0.76}}%
\end{pgfscope}%
\begin{pgfscope}%
\definecolor{textcolor}{rgb}{0.000000,0.000000,0.000000}%
\pgfsetstrokecolor{textcolor}%
\pgfsetfillcolor{textcolor}%
\pgftext[x=3.068386in,y=1.805406in,,bottom]{\color{textcolor}{\sffamily\fontsize{9.000000}{10.800000}\selectfont\catcode`\^=\active\def^{\ifmmode\sp\else\^{}\fi}\catcode`\%=\active\def%{\%}0.71}}%
\end{pgfscope}%
\begin{pgfscope}%
\definecolor{textcolor}{rgb}{0.000000,0.000000,0.000000}%
\pgfsetstrokecolor{textcolor}%
\pgfsetfillcolor{textcolor}%
\pgftext[x=5.009524in,y=1.965850in,,bottom]{\color{textcolor}{\sffamily\fontsize{9.000000}{10.800000}\selectfont\catcode`\^=\active\def^{\ifmmode\sp\else\^{}\fi}\catcode`\%=\active\def%{\%}0.79}}%
\end{pgfscope}%
\begin{pgfscope}%
\definecolor{textcolor}{rgb}{0.000000,0.000000,0.000000}%
\pgfsetstrokecolor{textcolor}%
\pgfsetfillcolor{textcolor}%
\pgftext[x=6.950661in,y=2.049783in,,bottom]{\color{textcolor}{\sffamily\fontsize{9.000000}{10.800000}\selectfont\catcode`\^=\active\def^{\ifmmode\sp\else\^{}\fi}\catcode`\%=\active\def%{\%}0.83}}%
\end{pgfscope}%
\begin{pgfscope}%
\definecolor{textcolor}{rgb}{0.000000,0.000000,0.000000}%
\pgfsetstrokecolor{textcolor}%
\pgfsetfillcolor{textcolor}%
\pgftext[x=1.476653in,y=1.985698in,,bottom]{\color{textcolor}{\sffamily\fontsize{9.000000}{10.800000}\selectfont\catcode`\^=\active\def^{\ifmmode\sp\else\^{}\fi}\catcode`\%=\active\def%{\%}0.80}}%
\end{pgfscope}%
\begin{pgfscope}%
\definecolor{textcolor}{rgb}{0.000000,0.000000,0.000000}%
\pgfsetstrokecolor{textcolor}%
\pgfsetfillcolor{textcolor}%
\pgftext[x=3.417791in,y=2.110297in,,bottom]{\color{textcolor}{\sffamily\fontsize{9.000000}{10.800000}\selectfont\catcode`\^=\active\def^{\ifmmode\sp\else\^{}\fi}\catcode`\%=\active\def%{\%}0.86}}%
\end{pgfscope}%
\begin{pgfscope}%
\definecolor{textcolor}{rgb}{0.000000,0.000000,0.000000}%
\pgfsetstrokecolor{textcolor}%
\pgfsetfillcolor{textcolor}%
\pgftext[x=5.358928in,y=1.802141in,,bottom]{\color{textcolor}{\sffamily\fontsize{9.000000}{10.800000}\selectfont\catcode`\^=\active\def^{\ifmmode\sp\else\^{}\fi}\catcode`\%=\active\def%{\%}0.71}}%
\end{pgfscope}%
\begin{pgfscope}%
\definecolor{textcolor}{rgb}{0.000000,0.000000,0.000000}%
\pgfsetstrokecolor{textcolor}%
\pgfsetfillcolor{textcolor}%
\pgftext[x=7.300066in,y=2.154608in,,bottom]{\color{textcolor}{\sffamily\fontsize{9.000000}{10.800000}\selectfont\catcode`\^=\active\def^{\ifmmode\sp\else\^{}\fi}\catcode`\%=\active\def%{\%}0.88}}%
\end{pgfscope}%
\begin{pgfscope}%
\pgfsetbuttcap%
\pgfsetmiterjoin%
\definecolor{currentfill}{rgb}{0.560784,0.658824,0.749020}%
\pgfsetfillcolor{currentfill}%
\pgfsetlinewidth{0.803000pt}%
\definecolor{currentstroke}{rgb}{0.435294,0.529412,0.611765}%
\pgfsetstrokecolor{currentstroke}%
\pgfsetdash{}{0pt}%
\pgfpathmoveto{\pgfqpoint{2.874519in}{2.417424in}}%
\pgfpathlineto{\pgfqpoint{3.152297in}{2.417424in}}%
\pgfpathlineto{\pgfqpoint{3.152297in}{2.514646in}}%
\pgfpathlineto{\pgfqpoint{2.874519in}{2.514646in}}%
\pgfpathlineto{\pgfqpoint{2.874519in}{2.417424in}}%
\pgfpathclose%
\pgfusepath{stroke,fill}%
\end{pgfscope}%
\begin{pgfscope}%
\definecolor{textcolor}{rgb}{0.000000,0.000000,0.000000}%
\pgfsetstrokecolor{textcolor}%
\pgfsetfillcolor{textcolor}%
\pgftext[x=3.263408in,y=2.417424in,left,base]{\color{textcolor}{\sffamily\fontsize{10.000000}{12.000000}\selectfont\catcode`\^=\active\def^{\ifmmode\sp\else\^{}\fi}\catcode`\%=\active\def%{\%}Sin Feature}}%
\end{pgfscope}%
\begin{pgfscope}%
\pgfsetbuttcap%
\pgfsetmiterjoin%
\definecolor{currentfill}{rgb}{0.850980,0.549020,0.584314}%
\pgfsetfillcolor{currentfill}%
\pgfsetlinewidth{0.803000pt}%
\definecolor{currentstroke}{rgb}{0.709804,0.431373,0.470588}%
\pgfsetstrokecolor{currentstroke}%
\pgfsetdash{}{0pt}%
\pgfpathmoveto{\pgfqpoint{4.324945in}{2.417424in}}%
\pgfpathlineto{\pgfqpoint{4.602722in}{2.417424in}}%
\pgfpathlineto{\pgfqpoint{4.602722in}{2.514646in}}%
\pgfpathlineto{\pgfqpoint{4.324945in}{2.514646in}}%
\pgfpathlineto{\pgfqpoint{4.324945in}{2.417424in}}%
\pgfpathclose%
\pgfusepath{stroke,fill}%
\end{pgfscope}%
\begin{pgfscope}%
\definecolor{textcolor}{rgb}{0.000000,0.000000,0.000000}%
\pgfsetstrokecolor{textcolor}%
\pgfsetfillcolor{textcolor}%
\pgftext[x=4.713834in,y=2.417424in,left,base]{\color{textcolor}{\sffamily\fontsize{10.000000}{12.000000}\selectfont\catcode`\^=\active\def^{\ifmmode\sp\else\^{}\fi}\catcode`\%=\active\def%{\%}Con Feature}}%
\end{pgfscope}%
\end{pgfpicture}%
\makeatother%
\endgroup%
}
    \caption{Average Precision con y sin ingeniería de variables, calculada en validación.}
    \label{fig:ap}
\end{figure}


El mejor desempeño de CP@20 y AP en validación correspondió a la red neuronal con ingeniería de variables. Frente a la regresión logística con variables históricas, su CP@20 mejoró 0,0152 puntos y su AP 0,0906 puntos. Sobre esta base se seleccionó la red neuronal para la única evaluación final en prueba.

\subsection{Evaluación final sobre el conjunto de prueba}
El conjunto de prueba filtrado comprendió 58.305 operaciones entre los días 157 y 179, de las cuales 113 fueron fraudulentas. La red neuronal seleccionada produjo los siguientes resultados:

\begin{table}[!htbp]
\centering
\caption{Resultados del modelo seleccionado en prueba}
\label{tab:test}
\begin{tabular}{@{}lrr@{}}
\toprule
\textbf{Métrica} & \textbf{Validación} & \textbf{Prueba} \\
\midrule
Card Precision@20 & 0,3000 & 0,2217 \\
Average Precision & 0,8890 & 0,8902 \\
ROC-AUC & 0,9993 & 0,9993 \\
\bottomrule
\end{tabular}
\par\smallskip
\end{table}

La disminución de CP@20 en prueba debe interpretarse considerando el número de clientes comprometidos disponibles. En validación, un total de 142 clientes distintos con fraude establece un máximo de $142/(23\times20)=0,3087$. El modelo alcanzó $138/(23\times20)=0,3000$. En prueba, los 104 clientes comprometidos fijan un máximo de $104/(23\times20)=0,2261$ y el resultado obtenido equivale a $102/(23\times20)=0,2217$. De este modo, aunque la métrica absoluta es menor en prueba, también lo es el máximo posible bajo el escenario evaluado.

En esta simulación se identificaron 102 de 104 clientes comprometidos presentes en prueba (98,1\%). Esta cifra no equivale a un \textit{recall} de transacciones ni garantiza la misma proporción de detección sobre datos reales. Del mismo modo, los valores elevados de AP y ROC-AUC se interpretan en el contexto del conjunto sintético y no constituyen evidencia suficiente de una generalización a una entidad bancaria.

\section{Impacto en el negocio}

La aplicación prevista consiste en ordenar diariamente las operaciones de mayor riesgo y agruparlas por cliente, a fin de presentar al equipo de prevención una lista priorizada de casos. En un escenario con capacidad para revisar 20 clientes por día, el procedimiento permite que la cantidad de alertas dependa de una restricción operativa explícita en lugar de un umbral arbitrario de probabilidad.

Durante los 23 días de prueba se consideran hasta 460 revisiones cliente-día ($20\times23$). Con la evaluación realizada, 102 selecciones corresponden a clientes con fraude y las restantes 358 no presentan fraude en el día en que fueron seleccionadas. Es importante observar que estas revisiones no equivalen necesariamente a 358 personas diferentes, porque un mismo cliente sin fraude puede volver a figurar en jornadas posteriores.

El resultado permite ilustrar un beneficio operativo potencial: priorizar casos relevantes en lugar de revisar operaciones sin orden de riesgo. Sin embargo, el dataset no contiene los costos reales de investigación, los importes recuperables ni la efectividad de las medidas posteriores al alerta. Por lo tanto, no es posible cuantificar de manera verificable un ahorro monetario o una reducción efectiva de pérdidas a partir de estos experimentos.

\subsection{Riesgos, aspectos éticos y limitaciones}
El carácter sintético de BankSim constituye una limitación central. Los patrones reproducidos por el simulador pueden diferir de los de una entidad real, y un desempeño elevado sobre estos datos no descarta problemas de generalización, cambios temporales de comportamiento o errores en nuevos escenarios.

Adicionalmente, la etiqueta de fraude se conoce con retraso, el dataset no distingue tarjetas independientes de un mismo cliente y su resolución temporal es diaria. Estos aspectos simplifican el problema operativo y deben revisarse antes de implementar un sistema de detección en tiempo real. También es necesario examinar la calibración de los puntajes: un buen ordenamiento de casos no implica automáticamente que los valores predichos puedan interpretarse como probabilidades fiables.

Los datos utilizados son sintéticos y no corresponden a transacciones personales identificables del mundo real \cite{inproceedings}. No obstante, una aplicación productiva requeriría controles de acceso, protección de datos personales, auditoría de las decisiones y revisión de posibles efectos adversos sobre grupos de clientes. La exclusión de edad y género del entrenamiento no demuestra por sí sola la ausencia de sesgos.

\subsection{Pasos futuros}
Como continuación del trabajo se propone evaluar otros períodos temporales, analizar la estabilidad de los resultados frente a cambios en la tasa de fraude y considerar una búsqueda controlada de hiperparámetros. También sería conveniente incorporar variables de anomalía de importe que utilicen exclusivamente antecedentes disponibles, revisar el tratamiento de operaciones del mismo día y comparar estrategias de calibración de puntajes.

Para avanzar hacia un caso empresarial real, sería necesario disponer de datos de operaciones y revisiones efectivas, estimar el costo de falsas alarmas y fraudes no detectados, medir tiempos de respuesta y contrastar los modelos con procesos de revisión existentes. Estos pasos son propuestas de trabajo y no forman parte de los experimentos realizados.

\section{Conclusiones y recomendaciones}

El proyecto permitió desarrollar un procedimiento de detección de fraude basado en aprendizaje supervisado, integrando exploración de datos, ingeniería de variables, validación temporal, comparación de modelos y una métrica orientada a la capacidad de investigación.

La evaluación mostró que las características históricas pueden aportar información útil, aunque su efecto depende del algoritmo utilizado. Entre las ocho alternativas comparadas, la red neuronal \textit{Feed-Forward} con feature engineering fue la seleccionada por su CP@20 y AP en validación. En el conjunto de prueba obtuvo CP@20 de 0,2217, AP de 0,8902 y ROC-AUC de 0,9993.

El principal aprendizaje metodológico es la importancia de definir indicadores coherentes con el problema de negocio. En este caso, la cantidad de clientes con fraude disponibles por día limita el valor alcanzable de CP@$k$; por eso, las diferencias entre modelos y períodos deben analizarse junto con esa restricción. Se recomienda conservar el enfoque temporal y el control de disponibilidad de etiquetas, y validar el sistema con datos más representativos antes de considerar cualquier uso operativo.

\begin{nota}
    Repositorio: \url{https://github.com/GinoHumberto/Fraud_detection}
\end{nota}

% --- BIBLIOGRAFÍA --- %
\printbibliography[title={Fuentes y Bibliografía}]

\end{document}